\documentclass[11pt,a4paper]{article}

\usepackage[italian]{babel}
\usepackage[utf8]{inputenc}
\usepackage[T1]{fontenc}
\usepackage[margin=3cm]{geometry}
\usepackage{amsmath}
\usepackage{graphicx}
\usepackage{caption}
\usepackage{setspace}
\usepackage{titlesec}
\usepackage{hyperref}
\usepackage{enumitem}
\usepackage{xcolor}
\usepackage{listings}

\renewcommand{\lstlistingname}{Codice}
\lstset{
  basicstyle=\ttfamily\small,
  frame=single,
  breaklines=true,
  columns=fullflexible,
  captionpos=b,
  backgroundcolor=\color{gray!5},
  numbers=none,
  showstringspaces=false
}

\hypersetup{
  colorlinks=true,
  linkcolor=black,
  urlcolor=blue,
  citecolor=black
}

\titleformat{\section}{\normalfont\Large\bfseries}{\thesection}{1em}{}
\titleformat{\subsection}{\normalfont\large\bfseries}{\thesubsection}{1em}{}

\setstretch{1.15}

\newcommand{\screenshotplaceholder}[2]{%
  \begin{figure}[h!]
    \centering
    \fbox{%
      \begin{minipage}[c][#1][c]{0.82\textwidth}
        \centering
        \color{gray}
        \textit{[spazio riservato allo screenshot: #2]}
      \end{minipage}%
    }
    \caption{#2}
  \end{figure}
}

\title{\textbf{Aspetti funzionali e implementativi dell'indice epigrafico digitale ILA}}
\author{}
\date{}

\begin{document}

\maketitle
\vspace{-2.5em}

\begin{center}
\textit{Bozza di lavoro --- versione interna}
\end{center}

\vspace{1.5em}

\begin{abstract}
\noindent
In questo lavoro vengono descritte le principali attività di sviluppo del
progetto ILA (\textit{Index Lunae Antiquae}), con particolare riferimento
agli aspetti realizzativi e funzionali dei risultati conseguiti. ILA si
propone di raccogliere, codificare e rendere fruibile con un'unica
interfaccia digitale l'insieme delle attestazioni epigrafiche relative al
culto del dio lunare frigio Men. Il progetto nasce dall'esigenza di
superare i limiti di un'edizione a stampa di riferimento --- il
\textit{Corpus Monumentorum Religionis Dei Menis} di E.\,N.\ Lane
(1971--1978) --- offrendo una rappresentazione digitale, interrogabile e
progressivamente rivedibile del medesimo materiale. La prima fase del
lavoro è stata dedicata al riconoscimento ottico dei testi a partire
dall'edizione a stampa e alla loro conversione in un formato di codifica
aperto e standardizzato. Nella fase successiva ci si è concentrati
sull'adozione e l'adattamento di un sistema di marcatura filologica
consolidato, sulla realizzazione di un'interfaccia di consultazione
multicanale e sulla progettazione di una sezione dedicata specificamente
alla descrizione iconografica dei monumenti, finora trattata in maniera
non sistematica nella letteratura di riferimento. Si presentano inoltre
le modalità di organizzazione dell'interfaccia, gli strumenti di ricerca
disponibili e le prospettive di sviluppo futuro del progetto.

\medskip
\noindent
This work describes the main development activities of the ILA
(\textit{Index Lunae Antiquae}) project, with particular reference to the
implementation and functional aspects of the results achieved. ILA aims
to collect, encode and make available within a single digital interface
the whole set of epigraphic attestations related to the cult of the
Phrygian lunar god Men. The project addresses the limits of a printed
reference edition by providing a digital, searchable and progressively
revisable representation of the same material, combining optical
character recognition of the source texts, a philological markup system,
and a dedicated iconographic description module.
\end{abstract}

\vspace{1em}

\section{Introduzione}

Il progetto ILA (\textit{Index Lunae Antiquae}) nasce con l'obiettivo di
raccogliere in un unico ambiente digitale l'intero corpus delle
attestazioni epigrafiche riconducibili al culto del dio lunare frigio
Men, un culto diffuso in età ellenistica e romana in vaste aree
dell'Anatolia occidentale e, per attestazioni più sporadiche, in altre
regioni del Mediterraneo antico. Il punto di partenza del lavoro è
rappresentato dall'edizione critica a stampa curata da E.\,N.\ Lane, il
\textit{Corpus Monumentorum Religionis Dei Menis} (CMRDM), pubblicata in
più volumi tra il 1971 e il 1978, che a tutt'oggi costituisce il
principale strumento di riferimento per lo studio sistematico delle
attestazioni relative a questa divinità.

Per quanto fondamentale, un'opera a stampa di questo tipo presenta i
limiti tipici del proprio formato: l'assenza di un indice interrogabile
per attributo (materiale, luogo di rinvenimento, tipologia di supporto,
epiteto divino, iconografia), l'impossibilità di aggiornare le singole
voci alla luce di revisioni successive senza una nuova edizione a
stampa, e una fruizione sostanzialmente lineare del contenuto, poco
adatta alle esigenze di uno studio comparativo su larga scala. Il
progetto ILA nasce proprio per superare questi limiti, proponendo una
rappresentazione digitale del medesimo patrimonio epigrafico, codificata
secondo standard aperti e internazionalmente riconosciuti, interrogabile
per singolo attributo e progressivamente rivedibile sulla base di un
confronto sistematico con le fonti primarie.

Il lavoro si è articolato in più fasi, non rigidamente sequenziali ma
in buona parte sovrapposte e tuttora in corso. Una prima fase è stata
dedicata al riconoscimento ottico dei testi (\textit{OCR}) a partire
dall'edizione a stampa di riferimento, propedeutico alla loro conversione
in formato digitale trattabile. Una seconda fase si è concentrata
sull'adozione di un sistema di codifica filologica adeguato alla
complessità del materiale, mutuandolo --- con i necessari adattamenti ---
da un sistema di marcatura già ampiamente collaudato in ambito
epigrafico digitale. Una terza fase, tuttora in corso, riguarda la
revisione sistematica delle singole voci sulla base di un confronto
diretto con le fonti primarie disponibili, così da ridurre
progressivamente la dipendenza dall'edizione a stampa originaria. In
parallelo a queste attività di codifica del testo, si è ritenuto
opportuno affiancare una sezione dedicata alla descrizione
dell'iconografia dei monumenti, aspetto che nella letteratura di
riferimento viene generalmente trattato solo in forma discorsiva,
all'interno del commento a ciascuna voce, e mai in una forma
sistematicamente interrogabile.

L'articolo è organizzato come segue. Nella sezione 2 viene illustrata
la modalità di rappresentazione concettuale e formale delle epigrafi.
Nella sezione 3 si descrive il sistema di codifica filologica adottato,
mutuato da un sistema di marcatura consolidato e adattato alle esigenze
specifiche del corpus, incluso il trattamento dei nomi divini e dei
rispettivi epiteti. La sezione 4 descrive la procedura di riconoscimento
ottico e trascrizione dei testi epigrafici. La sezione 5 presenta
l'organizzazione dell'interfaccia digitale, con particolare attenzione,
nella sottosezione 5.1, alla sezione dedicata all'iconografia dei
monumenti, nella sottosezione 5.2 alla navigazione tra divinità, epiteti
e registro onomastico, nella sottosezione 5.3 alla raccolta delle fonti
letterarie relative alla divinità e al lessico cultuale che le accomuna
al materiale epigrafico, e nella sottosezione 5.4 agli strumenti di
ricerca. Infine, nella sezione 6 si propongono alcune considerazioni
finali e le prospettive di sviluppo futuro del progetto.

\screenshotplaceholder{7cm}{Pagina iniziale dell'interfaccia digitale di ILA}

\section{Rappresentazione delle epigrafi}

Concettualmente, ogni epigrafe è concepita, in analogia con altri
progetti di museologia digitale, come una singola entità astratta
costituita da un insieme di attributi che ne caratterizzano il contenuto
e la forma materiale: materiale del supporto, tipologia del monumento,
luogo di rinvenimento e di conservazione, lingua dell'iscrizione, testo
dell'iscrizione, datazione, epiteto o epiteti del dio Men eventualmente
menzionati, e --- come si dirà più diffusamente nella sezione 5.1 ---
un insieme strutturato di descrittori iconografici.

A differenza di soluzioni che adottano un formato intermedio semplificato
per la gestione quotidiana dei dati, salvo poi generare da questo un
export in formato di codifica filologica standard, nel caso di ILA si è
scelto di adottare fin da subito la codifica filologica come formato
nativo e unico delle epigrafi. Questa scelta, più onerosa nella fase di
sviluppo iniziale, evita la comparsa di due rappresentazioni parallele
del medesimo dato --- una \virg{di comodo} per l'applicazione e
una \virg{scientifica} per la comunità di studiosi --- che nel
tempo rischierebbero di divergere silenziosamente l'una dall'altra. Ogni
epigrafe corrisponde così a un unico documento strutturato, che
costituisce al tempo stesso il dato di lavoro dell'applicazione e
l'edizione digitale consultabile e scaricabile dall'utente finale.

A partire da questa rappresentazione, un apposito modulo software
genera automaticamente le schede di consultazione presentate
nell'interfaccia grafica, così da garantire che la scheda visualizzata
dall'utente sia sempre generata direttamente dalla codifica filologica
sottostante, e non da una rappresentazione parallela soggetta a
disallineamento.

Lo stesso criterio guida anche gli interventi più recenti di
affinamento della rappresentazione, rivolti in particolare a due
attributi che in una prima stesura della codifica restavano mescolati
a un commento discorsivo più ampio. Il primo riguarda le note
paleografiche associate a ciascuna epigrafe, che in origine
contenevano, insieme al commento sulle forme delle lettere, sulle
legature o sull'interpunzione, anche la misura dell'altezza dei
caratteri, ereditata nella sua formulazione originaria dall'edizione a
stampa di riferimento e per questo spesso espressa in inglese
all'interno di un testo altrimenti redatto in italiano. Si è scelto di
isolare questa misura in un attributo strutturato a sé, distinto dal
commento paleografico vero e proprio, cosicché la scheda possa
mostrare separatamente un dato quantitativo, confrontabile tra voci
diverse, e un giudizio discorsivo sulla mano del lapicida, che tale
non è. Le schede compilate secondo la vecchia convenzione restano
leggibili senza reintervento manuale: la misura viene riconosciuta
all'interno del testo esistente, tanto nella sua formulazione italiana
quanto in quella originaria inglese, e ricondotta automaticamente al
nuovo attributo, cosicché la migrazione dell'intero corpus a questa
rappresentazione più puntuale non ha richiesto la revisione manuale di
ciascuna voce.

Un affinamento dello stesso tipo ha interessato l'apparato critico, la
nota che segnala, per una determinata porzione di testo, letture
alternative proposte da altri editori. Anche in questo caso
l'informazione era in origine un unico blocco di prosa (\virg{riga
cinque: ια secondo Lane}), che mescolava indistintamente tre dati di
natura diversa --- la collocazione della lettura all'interno del
testo, la lettura stessa e la fonte a cui va attribuita. Si è scelto
di scomporre la nota nei tre componenti corrispondenti, presentati
nella scheda in colonne distinte anziché come un'unica riga di prosa,
in modo che ciascuna informazione resti autonomamente leggibile e, in
prospettiva, interrogabile alla stregua degli altri attributi
epigrafici. Anche in questo caso le voci storiche, prive di
un'attribuzione esplicita della fonte, vengono scomposte
automaticamente al momento della lettura, isolando la sequenza di nomi
propri che tipicamente chiude la nota dal resto del commento, così da
non richiedere una revisione sistematica del corpus pregresso.

\screenshotplaceholder{8cm}{Esempio di scheda epigrafica nell'interfaccia di consultazione}

\section{Il sistema di marcatura filologica}

Per la codifica delle epigrafi si è scelto di non progettare un sistema
di marcatura ex novo, bensì di mutuare --- adattandolo alle specifiche
esigenze del corpus --- un sistema di marcatura filologica già
largamente adottato e collaudato nell'ambito dell'epigrafia digitale,
basato sulla notazione convenzionale nota come \textit{sistema di
Leida}, tradizionalmente impiegata dagli epigrafisti per rappresentare
in forma sintetica i fenomeni testuali riscontrabili sui supporti
antichi (lacune, integrazioni, lettere dubbie, cancellature, restauri
editoriali, e così via).

Questa scelta ha comportato un lavoro di adattamento su due livelli.
Da un lato, è stato necessario costruire una corrispondenza puntuale e
sistematica tra ciascuna convenzione della notazione di Leida e il
corrispondente costrutto di marcatura nel formato di codifica digitale
adottato per il progetto, in modo che ogni fenomeno testuale --- una
lacuna, un'integrazione, una lettera di lettura incerta --- trovasse una
rappresentazione univoca e non ambigua. Dall'altro lato, si è reso
necessario introdurre uno strumento di supporto alla codifica che
permettesse di applicare queste convenzioni in maniera assistita,
riducendo il rischio di errore nella trascrizione dei fenomeni testuali
più complessi e rendendo il processo di marcatura accessibile anche a chi
non abbia una lunga esperienza diretta con la sintassi del formato di
codifica sottostante.

Un caso specifico, emerso nella pratica quotidiana di trascrizione,
riguarda la combinazione tra una lacuna materiale del supporto e
un'integrazione editoriale che occupa, in tutto o in parte, la
medesima porzione di testo. La convenzione a stampa rende normalmente
questo fenomeno racchiudendo l'intera sequenza --- lacuna e
integrazione insieme --- in un'unica coppia di parentesi quadre; una
prima resa nell'interfaccia trattava invece i due fenomeni come
indipendenti, producendo due coppie di parentesi annidate l'una
nell'altra anziché una sola. Si è pertanto introdotta un'azione di
marcatura dedicata a questo caso, che riconosce quando un'integrazione
racchiude direttamente una lacuna adiacente e restituisce, coerentemente
con la convenzione a stampa, un'unica coppia di parentesi per l'intera
sequenza, evitando una doppia parentesizzazione che, per quanto
formalmente equivalente, risulterebbe fuorviante per un lettore
abituato alla notazione tradizionale.

A questo sistema di marcatura mutuato dalla tradizione epigrafica si
affiancano, in analogia con le proposte già avanzate in altri progetti
di museologia epigrafica digitale, alcune estensioni pensate per
esplicitare informazioni che nei corpora epigrafici tradizionali
restano tipicamente implicite nel commento discorsivo: l'ambito
religioso dell'iscrizione, la collezione o il contesto di provenienza,
lo stato di conservazione e --- aspetto su cui si tornerà nella sezione
5.1 --- la descrizione strutturata dell'iconografia figurata associata
al monumento.

Un caso che ha richiesto particolare cura è la codifica dei nomi divini
e dei rispettivi epiteti cultuali. Il dio Men compare nel corpus non di
rado insieme ad altre divinità, in formule sincretiche o giustapposte
(un dio solare, una grande madre, una coppia di divinità associate nel
medesimo santuario), e ciascuna di queste figure può portare i propri
epiteti distinti. Trattare l'insieme degli epiteti di un'iscrizione come
una lista piatta, comune a tutte le divinità nominate nel testo,
produrrebbe un risultato filologicamente scorretto: attribuirebbe a un
dio un epiteto che il testo riferisce in realtà a un altro, semplicemente
perché i due nomi compaiono nella stessa iscrizione. Si è scelto
pertanto di codificare l'associazione tra ciascuna divinità e i propri
epiteti a livello della singola menzione all'interno del testo, cosa che
richiede, in fase di marcatura, di riconoscere correttamente dove finisce
il nome della divinità (che può essere esso stesso composto da più
parole, come nel caso di formule sincretiche) e dove comincia l'epiteto
che lo qualifica --- una distinzione non sempre ovvia nella resa a stampa
dell'edizione di riferimento, e che in una prima stesura della codifica
non era stata trattata con la necessaria sistematicità.

L'applicazione di questo criterio è stata condotta in maniera sistematica
su tutte le voci del corpus, verificando e correggendo, epigrafe per
epigrafe, sia la segmentazione tra nome divino ed epiteto sia i casi in
cui una classificazione precedente risultava imprecisa. Si tratta di un
lavoro di revisione capillare, più vicino a un controllo di qualità
filologico che a uno sviluppo di funzionalità in senso stretto, ma che
condiziona direttamente l'affidabilità di qualunque strumento --- di
consultazione, di ricerca o di analisi comparativa --- costruito a valle
su questa informazione (si veda, in particolare, la sezione 5.2).

Una parte non trascurabile di questo lavoro di adattamento ha riguardato
la conversione, su scala dell'intero corpus, delle convenzioni di Leida
ancora presenti come semplice punteggiatura all'interno del testo
trascritto --- parentesi quadre, parentesi tonde, sequenze di puntini ---
nei corrispondenti elementi di marcatura strutturata (integrazioni
editoriali, lacune, letture incerte, scioglimenti di abbreviazione). Fino
a quel momento, infatti, buona parte delle epigrafi del corpus recava sì
la notazione convenzionale corretta, ma non ancora tradotta nel formato
di codifica: una situazione che, pur non compromettendo la leggibilità
per un occhio esperto, rendeva quell'informazione invisibile a qualunque
strumento automatico costruito sulla marcatura strutturata, dagli stessi
meccanismi di ricerca alla generazione delle schede di consultazione. Si
sono pertanto realizzati due strumenti dedicati a questa conversione,
applicati sistematicamente a tutte le voci del corpus: il primo distingue,
all'interno delle parentesi quadre, il testo materialmente perduto dal
supporto dai puntini che segnalano invece lettere semplicemente non più
leggibili su un supporto altrimenti integro --- una distinzione consolidata
nella pratica epigrafica ma facile da perdere quando resta affidata alla
sola punteggiatura; il secondo distingue, all'interno delle parentesi
tonde, lo scioglimento di un'abbreviazione dalle lettere omesse per
errore dal lapicida, verificando caso per caso, sull'edizione a stampa di
riferimento, i casi in cui questa distinzione non fosse ricavabile dal
solo contesto immediato. In entrambi i casi si è seguito lo stesso
principio metodologico già enunciato nella sezione 4: la quantità di
testo mancante non viene mai dedotta dal numero di puntini presenti nella
trascrizione, poiché tale conteggio non è di per sé affidabile, ma viene
conservata separatamente per un'eventuale verifica successiva
sull'edizione a stampa.

Ad affiancare questi strumenti di conversione è stato introdotto anche un
dispositivo di controllo automatico che verifica, sull'intero corpus e
non voce per voce, la coerenza interna della marcatura: il bilanciamento
delle parentesi, la corrispondenza tra un valore numerico marcato
esplicitamente e il numerale greco effettivamente trascritto, e la
presenza di un elemento di datazione strutturata laddove il testo
dell'iscrizione riporta esplicitamente una formula d'anno. Un simile
controllo, eseguito sistematicamente anziché a campione, ha permesso di
individuare e correggere alcuni casi in cui l'origine dell'incoerenza non
era la marcatura in sé ma un errore di trascrizione a monte, confermando
quanto già osservato a proposito della revisione capillare descritta più
sopra: il controllo automatico della coerenza strutturale non sostituisce
la verifica filologica sulle fonti primarie, ma la orienta, segnalando
sistematicamente dove condurla in via prioritaria.

Lo stesso lavoro di affinamento ha infine riguardato lo strumento
assistito di marcatura in sé: in una prima versione, le azioni per
inserire una lacuna o uno spazio vacante nel testo potevano operare
soltanto in coda al testo esistente oppure in sostituzione di una
porzione di testo già selezionata, rendendo ad esempio impossibile
segnalare una lacuna a inizio riga senza cancellare i caratteri che
seguivano. Lo strumento consente ora di fissare un punto di inserimento
con un semplice clic all'interno del testo, e di inserire lì la lacuna o
lo spazio vacante senza alterare il contenuto circostante --- un
accorgimento minore nella sua descrizione, ma che rimuove un vincolo che
nella pratica quotidiana della marcatura si era rivelato più limitante
del previsto.

\screenshotplaceholder{7.5cm}{Strumento assistito di marcatura filologica in azione su una voce del corpus}

\section{Trascrizione e riconoscimento ottico dei testi}

Particolare attenzione è stata dedicata alla trascrizione dei testi
epigrafici a partire dall'edizione a stampa di riferimento e alla loro
successiva codifica filologica. Nella maggior parte dei casi, i testi
da trascrivere sono stati prelevati dalle voci corrispondenti
dell'edizione critica a stampa attraverso un processo di estrazione
semiautomatizzato, basato sul riconoscimento ottico dei caratteri
(\textit{OCR}), applicato alla versione digitalizzata del volume di
riferimento.

Questa fase si è rivelata particolarmente delicata a causa della natura
del materiale: i testi greci ed anatolici presenti nell'edizione,
spesso corredati di segni diacritici, note editoriali sopralineari o
sottolineari e simboli non corrispondenti al set di caratteri latino di
base, non sono sempre stati riconosciuti correttamente dagli strumenti
di riconoscimento ottico impiegati. In particolare, i caratteri greci
accentati e le convenzioni tipografiche utilizzate nell'edizione per
segnalare restauri, lettere ambigue o cancellature testuali si sono
dimostrati i punti più critici del processo, rendendo necessaria una
fase sistematica di revisione manuale dell'output prodotto dal
riconoscimento automatico, prima di procedere alla codifica vera e
propria del testo secondo il sistema di marcatura descritto nella
sezione precedente.

Il testo ottenuto dal riconoscimento ottico, una volta corretto, viene
dunque arricchito manualmente dei fenomeni testuali non direttamente
desumibili dall'OCR (interruzioni di riga, restauri editoriali,
caratteri ambigui) e infine convertito nella marcatura filologica
descritta nella sezione 3. Nei casi in cui il testo non risultasse
adeguatamente ricostruibile a partire dalla sola edizione a stampa, si è
proceduto, ove possibile, a un confronto diretto con le fonti primarie
disponibili, secondo il principio metodologico per cui nessun dato
mancante viene stimato o integrato arbitrariamente: un'informazione non
attestata resta dichiaratamente assente, piuttosto che essere colmata
per congettura non verificabile.

Poiché la revisione di una singola voce può coinvolgere più interventi
distinti nel tempo --- trascrizione iniziale, correzione dell'OCR,
integrazione filologica, revisione finale --- e non necessariamente ad
opera della stessa persona, si è ritenuto utile affiancare allo
strumento di marcatura un piccolo pannello, sempre visibile durante
l'editing, che permette di annotare rapidamente chi ha svolto un
determinato intervento e quando, senza dover interrompere il lavoro di
codifica per raggiungere una sezione separata dedicata alla cronologia
delle revisioni. Anche in questo caso si tratta di un'informazione che
viene scritta in un campo già previsto dalla codifica filologica di
base, e non di un dato collaterale gestito a parte: la tracciabilità
delle singole mani che hanno contribuito a una voce resta così parte
integrante della scheda, anziché affidata alla memoria di chi coordina
il lavoro editoriale.

\screenshotplaceholder{6.5cm}{Confronto tra testo a stampa e trascrizione ottenuta tramite riconoscimento ottico}

Il testo digitale già incorporato nell'edizione a stampa di riferimento
pone un problema specifico quando si tratta di verificare, per una
singola voce, la lettura originaria del testo greco: il livello testuale
del documento digitalizzato rende correttamente le porzioni in caratteri
latini, ma restituisce le porzioni in greco come una sequenza di
caratteri privi di senso, esito di una corrispondenza tra il font
tipografico originale e la codifica dei caratteri che non riproduce un
alfabeto greco standard. Applicare indiscriminatamente alla pagina intera
uno strumento di riconoscimento ottico specializzato per il greco non
risolve il problema, perché un simile strumento tende a leggere come
greco anche le porzioni in inglese del commento circostante,
introducendo un rumore che complica il confronto più di quanto non lo
agevoli. Si è scelto pertanto un procedimento più mirato: a partire
dalla posizione di ciascuna riga sulla pagina, si distinguono le righe
effettivamente redatte in inglese da quelle che, pur apparendo come
sequenze illeggibili, corrispondono in realtà a testo greco mal
codificato --- riconoscibili da alcuni indizi strutturali, come la
presenza di cifre all'interno di una parola o di maiuscole in posizione
mediana, estranee alla tipografia inglese corrente --- per poi ritagliare
dalla sola immagine della pagina i blocchi corrispondenti alle righe
greche e sottoporli, isolatamente, al riconoscimento ottico specializzato.
La pagina ricomposta unisce così il testo inglese, già attendibile a
partire dal livello testuale originario, e il testo greco, ottenuto per
la sola porzione che lo richiede da un riconoscimento ottico mirato
anziché generico.

Il confronto sistematico tra la trascrizione presente nel corpus e la
lettura offerta da un'edizione a stampa, già indicato altrove in questo
lavoro come uno dei fronti di sviluppo in corso, ha trovato in questa
fase un primo strumento operativo dedicato. Una procedura di collazione
assistita affianca, voce per voce, il testo greco del corpus e quello
della corrispondente edizione a stampa presa a riferimento --- non
necessariamente il solo CMRDM, ma una qualunque delle edizioni a stampa
ormai riconosciute dal sistema secondo un registro esteso, che affianca
all'opera di Lane altre raccolte epigrafiche di riferimento ---
classificando ogni divergenza riscontrata per tipologia (una singola
lettera, un numerale, un intero blocco di testo) e raggruppando gli
scambi che ricorrono più di frequente, poiché è precisamente in queste
classi di errore sistematico, e non nelle singole discrepanze isolate,
che conviene concentrare la revisione filologica.

Una prima applicazione di questo strumento a un'intera edizione del
corpus di riferimento ha mostrato, tuttavia, che affidarsi a una sola
lettura automatica della fonte a stampa produce un numero di divergenze
poco significativo, perché la classe di errore più frequente non era
imputabile al corpus, ma allo stesso riconoscimento ottico: una
confusione sistematica e unidirezionale tra due lettere greche
graficamente simili, che nessun trascrittore umano commetterebbe con la
medesima regolarità. Si è scelto pertanto di raddoppiare la lettura
automatica della fonte a stampa, impiegando due motori di riconoscimento
ottico distinti, che si è osservato commettere errori sistematici
differenti l'uno dall'altro: una divergenza rispetto al corpus viene
riportata come tale soltanto quando entrambe le letture concordano nel
segnalarla, e la misura di somiglianza adottata è quella della lettura
più vicina al testo del corpus, cosicché, quando una delle due letture
automatiche coincide con il corpus, la discrepanza dell'altra viene
correttamente attribuita a un errore del riconoscimento ottico e non a
un problema della trascrizione. Questo doppio controllo incrociato
riduce sensibilmente il rumore introdotto dal solo processo di OCR,
lasciando emergere con maggiore nitidezza le divergenze che meritano
davvero una verifica sulle fonti primarie.

\section{Organizzazione dell'interfaccia digitale}

Al fine di rendere il patrimonio epigrafico raccolto fruibile sia agli
specialisti della materia sia a un pubblico più ampio di visitatori non
specialisti, l'interfaccia digitale è stata progettata secondo criteri
di \textit{responsive web design}, così da risultare consultabile
indipendentemente dal dispositivo utilizzato.

L'interfaccia si articola in diverse sezioni tra loro complementari:
una sezione di consultazione che presenta l'elenco completo delle
epigrafi con accesso diretto alle singole schede descrittive; una
sezione bibliografica che raccoglie i repertori e i corpora di
riferimento; una sezione geografica che consente di visualizzare la
distribuzione spaziale delle attestazioni su una mappa interattiva,
affiancata da una sezione cronologica che ne restituisce invece la
distribuzione lungo l'asse del tempo, disposta come una linea temporale
navigabile e progressivamente ingrandibile; una sezione dedicata alle
fonti letterarie sulla divinità, descritta nella sottosezione 5.3; e
una sezione di ricerca avanzata, descritta più diffusamente nella
sottosezione 5.4. Ogni scheda descrittiva è corredata, ove disponibile,
di materiale fotografico del monumento e riporta, oltre al testo
dell'iscrizione e ai principali attributi epigrafici, un rimando
puntuale alla codifica filologica sottostante, liberamente
consultabile e scaricabile.

A questo insieme di rimandi si affianca, per ciascuna epigrafe, un
ulteriore livello di collegamento verso i principali repertori
digitali di riferimento per l'epigrafia antica --- tra gli altri, le
banche dati che raccolgono trascrizioni e immagini su scala regionale
o tematica, e i grandi indici prosopografici e geografici del mondo
antico. Ogni riferimento è codificato come una coppia costituita dalla
sigla del repertorio e dall'identificativo assegnato dalla risorsa
esterna, da cui l'interfaccia ricava --- per i repertori più diffusi
--- il collegamento diretto, lasciando comunque la possibilità di
correggerlo a mano nei casi meno standard. Anche in questo caso si è
scelto un vocabolario di sigle aperto, anziché un elenco chiuso
stabilito una volta per tutte: le sigle già presenti nel corpus
costituiscono l'elenco di base proposto in fase di redazione, ma è
possibile introdurne di nuove senza dover intervenire sullo schema di
codifica, secondo lo stesso principio già illustrato a proposito del
vocabolario iconografico (sottosezione 5.1). Nella scheda, questi
riferimenti compaiono come elementi cliccabili accanto agli altri
identificativi bibliografici dell'epigrafe, così da rendere il
passaggio verso la risorsa esterna corrispondente immediato quanto la
consultazione di un rimando interno al corpus.

Fino a poco tempo fa, lo stato di consultazione dell'interfaccia viveva
soltanto nella memoria della sessione in corso: aprire una scheda,
ricaricare la pagina o condividerne l'indirizzo con un collega
significava perdere il punto in cui ci si trovava e ripartire dalla
pagina iniziale, senza alcun modo di indirizzare qualcuno direttamente
verso \emph{quella} scheda. Si è scelto pertanto di far viaggiare,
nell'indirizzo stesso della pagina, sia la sezione dell'interfaccia
attualmente aperta sia, quando pertinente, la singola scheda consultata,
così da rendere possibile un collegamento diretto che salti la pagina
iniziale e conduca chi lo riceve esattamente alla voce che si intendeva
condividere. Questa informazione è stata inserita come parte
interrogabile dell'indirizzo, anziché come segmento del percorso, poiché
l'infrastruttura su cui l'interfaccia è pubblicata non consente di
riscrivere le rotte lato server, e un collegamento di questo tipo deve
poter funzionare anche in un ambiente di pubblicazione puramente
statico.

Sotto il titolo di ciascuna scheda compaiono ora, alla portata di
chiunque consulti il corpus e non soltanto di chi dispone dei permessi
di modifica, tre elementi pensati per agevolare il riuso scientifico del
materiale: il collegamento diretto appena descritto, lo scarico del
documento filologico sottostante, già previsto fin dall'inizio del
progetto (si veda la sezione 2), e una proposta di citazione già pronta
all'uso. Quest'ultima elenca, nell'ordine in cui compaiono nel registro
delle edizioni a stampa riconosciute dal sistema --- esteso, come si è
detto, ben oltre la sola opera di Lane --- tutti i riferimenti
bibliografici a stampa individuati per quella scheda, componendo una
stringa del tipo \virg{ILA-042 (= CMRDM I 72 = Lane, Berytus 15, n.
38)}, che restituisce in un colpo d'occhio la concordanza tra la
numerazione interna del progetto e le diverse numerazioni con cui la
medesima epigrafe è già nota alla tradizione a stampa.

La visualizzazione cronologica applica alla dimensione temporale lo
stesso principio di trasparenza già adottato altrove nel progetto (si
veda la sezione 4): poiché non tutte le epigrafi del corpus recano
estremi cronologici numericamente determinabili, la linea del tempo
colloca soltanto le voci per cui questa informazione è disponibile, ma
dichiara sempre, accanto al conteggio delle voci effettivamente
collocate, quante schede ne restano escluse per assenza di datazione,
con la possibilità di consultarne l'elenco completo con un solo
passaggio: l'assenza di una scheda dalla linea del tempo deve poter
essere sempre distinta da un errore, e non lasciata all'interpretazione
di chi consulta.

\screenshotplaceholder{8cm}{Vista d'insieme dell'elenco delle epigrafi con filtri applicati}

\screenshotplaceholder{8cm}{Visualizzazione geografica delle attestazioni su mappa interattiva}

\screenshotplaceholder{7.5cm}{Visualizzazione cronologica delle attestazioni, con panoramica per secolo e conteggio delle schede prive di datazione}

\subsection{La sezione iconografica}

Uno degli aspetti a cui si è ritenuto di dedicare uno sviluppo autonomo
è la descrizione dell'iconografia figurata associata ai monumenti. Nella
letteratura di riferimento, l'apparato iconografico delle stele e dei
rilievi dedicati a Men viene generalmente descritto in forma discorsiva,
all'interno del commento a ciascuna voce, senza una struttura che ne
consenta l'interrogazione sistematica e il confronto comparativo tra
voci diverse. Si è pertanto scelto di introdurre, per ciascuna epigrafe
dotata di apparato figurato, una sezione dedicata alla descrizione
strutturata dell'iconografia, articolata secondo un vocabolario
controllato di categorie ricorrenti: la funzione compositiva della
scena, la tipologia e l'identificazione delle singole figure
rappresentate, e un insieme di tratti descrittivi specifici --- copricapo,
attributi lunari, oggetti tenuti in mano, animali associati, abbigliamento,
tecnica esecutiva e posizione compositiva della figura all'interno del
rilievo.

Questa strutturazione consente, da un lato, di visualizzare in maniera
sintetica e coerente l'apparato iconografico di ciascun monumento
all'interno della relativa scheda descrittiva, e dall'altro di rendere
tale informazione interrogabile alla stregua degli altri attributi
epigrafici, aprendo la possibilità di ricerche trasversali --- ad esempio
l'insieme di tutte le raffigurazioni in cui il dio è rappresentato a
cavallo, oppure con un determinato copricapo o oggetto in mano --- che
sarebbero difficilmente praticabili a partire dalla sola lettura
discorsiva dell'edizione a stampa.

Dal punto di vista della codifica, la scelta di dove collocare questa
informazione all'interno del documento filologico non era scontata.
L'apparato iconografico non corrisponde a nessuno degli elementi
previsti dallo schema di codifica di base per la descrizione testuale
dell'epigrafe, e forzarlo dentro elementi pensati per altri scopi
avrebbe reso la codifica ambigua e difficile da estendere in futuro. Si
è scelto pertanto di ospitarlo in un blocco di dati extra --- lo spazio
che lo standard di codifica riserva esplicitamente a informazioni non
previste dallo schema di base --- così da tenere nettamente separati il
testo dell'iscrizione, filologicamente vincolante, e la sua descrizione
iconografica, che resta per sua natura un'interpretazione soggetta a
revisione. Questa collocazione ha un vantaggio pratico non trascurabile:
il blocco può essere esteso con nuove categorie descrittive senza dover
intervenire sullo schema di codifica generale del corpus, ed è del tutto
opzionale, per cui una scheda priva di apparato figurato resta valida
senza dover \virg{inventare} una descrizione iconografica assente.

Il Codice~\ref{lst:xenodata} mostra, semplificato ma realistico, il
blocco corrispondente a un rilievo votivo in cui Men è raffigurato al
centro della composizione, con berretto frigio, falce lunare sulle
spalle e una pigna nella mano. Si noti come ogni informazione sia
espressa come coppia \emph{tipo}/\emph{valore} (l'attributo
\texttt{type} e l'attributo \texttt{key} di ciascun elemento
\texttt{trait}), anziché come testo libero: è proprio questa scelta,
apparentemente minore, a rendere l'informazione machine-readable senza
sacrificarne la leggibilità per un occhio umano che sappia già leggere
un documento TEI.

\begin{lstlisting}[language=XML, caption={Blocco \texttt{xenoData} di
esempio per un rilievo votivo con Men frontale}, label={lst:xenodata}]
<xenoData>
    <iconography>
        <function key="votive"/>
        <figure n="1" type="deity" key="Men" place="top_centre">
            <trait type="headgear" key="phrygian_cap"/>
            <trait type="lunar" key="crescent_shoulders"/>
            <trait type="held_object" key="pine_cone"/>
        </figure>
        <note>Rilievo votivo, Men frontale con falce lunare
            sulle spalle.</note>
    </iconography>
</xenoData>
\end{lstlisting}

Proprio perché la descrizione dei tratti iconografici è vincolata a un
vocabolario chiuso e non a testo libero, ogni valore inserito in una
scheda è per costruzione confrontabile con lo stesso valore inserito in
qualunque altra scheda del corpus: \virg{copricapo radiato} indica
sempre lo stesso concetto, indipendentemente da chi abbia compilato la
scheda o da come il tratto sia descritto nel commento discorsivo
dell'edizione a stampa. Questa uniformità è ciò che rende possibile un
secondo strumento di analisi, pensato non per la consultazione della
singola epigrafe ma per lo studio comparativo dell'intero corpus: una
heatmap delle co-occorrenze tra coppie di attributi (iconografici, ma
anche onomastici e geografici), che misura quanto frequentemente due
valori compaiono insieme nelle stesse schede rispetto a quanto ci si
attenderebbe se fossero distribuiti in maniera indipendente. Senza un
vocabolario standardizzato alle spalle, una simile analisi sarebbe
impraticabile: basterebbe una minima variazione nella formulazione di
uno stesso tratto --- \virg{in groppa a un cavallo} anziché \virg{a
cavallo} --- per far apparire come distinti due fenomeni che sono in
realtà identici, falsando qualunque conteggio di co-occorrenza. La
coerenza terminologica garantita dal vocabolario controllato è dunque
condizione necessaria, non solo per la leggibilità della singola scheda,
ma per l'affidabilità di ogni analisi che aggreghi i dati sull'intero
corpus.

Un vocabolario chiuso pone però un problema speculare a quello che
risolve: se le chiavi ammesse sono fissate una volta per tutte, chi
compila una scheda deve conoscerle a memoria o consultare costantemente
un riferimento esterno, il che è precisamente l'ostacolo che si voleva
evitare rendendo l'iconografia accessibile anche a chi non abbia una
formazione specialistica in epigrafia. La soluzione adottata separa
nettamente due livelli: da un lato la codifica, che lavora sempre e solo
con chiavi sintetiche in inglese (\texttt{phrygian\_cap},
\texttt{crescent\_shoulders}, \texttt{pine\_cone}, come nell'esempio del
Codice~\ref{lst:xenodata}); dall'altro un vocabolario di traduzione, che
associa a ciascuna chiave un'etichetta discorsiva in italiano, pensata
per essere presentata a chi consulta o compila una scheda senza che
questi debba mai vedere o scrivere direttamente la chiave sottostante.
Il Codice~\ref{lst:vocab} ne mostra un estratto minimale.

\begin{lstlisting}[caption={Estratto del vocabolario di traduzione
chiave $\to$ etichetta}, label={lst:vocab}]
phrygian_cap        ->  "berretto frigio"
crescent_shoulders  ->  "falce lunare sulle spalle"
pine_cone           ->  "pigna"
votive               ->  "votiva"
top_centre           ->  "al centro in cima"
\end{lstlisting}

Questa separazione ha una conseguenza pratica diretta sulla facilità
d'uso: chi compila una scheda sceglie da un elenco di etichette già in
italiano --- \virg{berretto frigio}, \virg{falce lunare sulle spalle},
\virg{pigna} --- senza dover conoscere né scrivere la sintassi
sottostante, che viene generata automaticamente a partire dalla scelta
fatta. Allo stesso tempo, estendere il vocabolario --- perché
un'iscrizione presenta un attributo iconografico non ancora catalogato,
come è di fatto accaduto più volte nel corso della revisione del corpus
--- si riduce ad aggiungere una singola coppia chiave/etichetta a questa
tabella, senza dover intervenire né sullo schema di codifica descritto
poco sopra, né sulla logica applicativa che lo consulta: la tabella è,
in questo senso, l'unico punto del sistema che deve cambiare quando
cambia il vocabolario, e i vantaggi di questa scelta valgono
esattamente lo stesso principio di \virg{unica fonte di verità} già
discusso a proposito della rappresentazione dei dati (sezione 2).

Vale la pena spiegare, in parole semplici, in che modo la heatmap
quantifica la forza di un'associazione, perché il numero mostrato in
ciascuna cella non è un semplice conteggio di co-occorrenze, ma una
misura che tiene conto anche di quanto ciascun attributo sia, di per sé,
frequente o raro nel corpus. Due attributi molto comuni --- una regione
molto rappresentata, un materiale molto diffuso --- compariranno spesso
insieme semplicemente per effetto della loro frequenza individuale,
senza che questo indichi un legame significativo tra i due; viceversa,
due attributi rari che si presentano ripetutamente insieme, anche solo
poche volte in termini assoluti, segnalano un'associazione che difficilmente
può essere frutto del caso. La misura impiegata, nota in letteratura come
\emph{informazione mutua puntuale} (\emph{pointwise mutual information},
PMI), è pensata esattamente per isolare questo secondo tipo di
segnale, confrontando la frequenza osservata della coppia con la
frequenza che ci si aspetterebbe se i due attributi fossero tra loro
del tutto indipendenti.

Indicando con $P(A)$ e $P(B)$ la proporzione di schede del corpus in cui
compare, rispettivamente, il valore $A$ e il valore $B$, e con $P(A,B)$
la proporzione di schede in cui i due valori compaiono insieme, la
misura è definita come

\[
  \mathrm{PMI}(A,B) \;=\; \log_2 \left( \frac{P(A,B)}{P(A)\cdot P(B)} \right).
\]

Il denominatore $P(A) \cdot P(B)$ rappresenta la frequenza di
co-occorrenza che si osserverebbe se $A$ e $B$ comparissero nelle schede
in maniera del tutto scorrelata l'uno dall'altro; il numeratore $P(A,B)$
è invece la frequenza di co-occorrenza effettivamente osservata nel
corpus. Il logaritmo in base 2 di questo rapporto ha una lettura
intuitiva: quando le due frequenze coincidono, il rapporto vale $1$ e il
logaritmo si annulla ($\mathrm{PMI}=0$), a indicare nessuna associazione
oltre il livello atteso per puro caso; quando i due attributi compaiono
insieme più spesso di quanto atteso, il rapporto supera $1$ e la PMI
assume valore positivo, tanto più alto quanto più l'associazione è
marcata; quando invece compaiono insieme meno spesso del previsto, il
rapporto scende sotto $1$ e la PMI diventa negativa. La scelta della base
$2$, anziché del logaritmo naturale, non altera questa lettura --- ne
cambia solo la scala --- ma ha il vantaggio di rendere il valore
interpretabile in \virg{raddoppi}: un incremento di un'unità di PMI
corrisponde a una frequenza di co-occorrenza osservata doppia rispetto a
quella attesa per indipendenza. Nell'interfaccia, questo valore non viene
mai presentato da solo, ma sempre accompagnato dal conteggio assoluto
delle co-occorrenze sottostanti: una PMI elevata calcolata su due o tre
casi resta un'indicazione da verificare sulle fonti, non una conclusione,
ed è segnalata come tale.

La PMI a due fattori descritta sopra risponde bene alla domanda
\virg{questi due attributi si associano più del previsto?}, ma non
cattura un tipo di segnale altrettanto interessante per uno studio del
culto: combinazioni di \emph{tre} o più caratteristiche che, prese a
coppie, non risultano particolarmente marcate, ma che insieme
individuano un pattern ricorrente --- ad esempio un determinato epiteto
che, in una specifica regione, compare quasi sempre su un preciso tipo
di supporto. Per intercettare anche questo tipo di associazione, alla
modalità a coppie si affianca una modalità alternativa, che combina fino
a tre fattori in una singola \virg{colonna composita} e ne confronta la
probabilità congiunta osservata con quella attesa nell'ipotesi che tutti
i fattori coinvolti siano reciprocamente indipendenti --- una
generalizzazione diretta della stessa idea a $n$ variabili anziché due:

\[
  \mathrm{PMI}(A_1, A_2, \ldots, A_n) \;=\; \log_2 \left(
  \frac{P(A_1, A_2, \ldots, A_n)}{P(A_1)\cdot P(A_2)\cdots P(A_n)}
  \right).
\]

La lettura resta la stessa della formula a due fattori --- zero
per l'indipendenza reciproca, positivo per un'associazione più forte del
previsto, negativo per una più debole --- semplicemente estesa a un
denominatore che è ora il prodotto di tutte le frequenze marginali
coinvolte anziché di due sole. Va da sé che, aumentando il numero di
fattori combinati, il numero di schede che condividono \emph{tutti} i
valori scelti si riduce rapidamente, e con esso l'affidabilità
statistica della stima: la stessa cautela già raccomandata per la PMI a
due fattori sui campioni piccoli si applica qui con ancora più forza, ed
è per questo che il conteggio assoluto resta sempre visibile accanto al
valore, qualunque sia il numero di fattori combinati.

Il motivo per cui questo intero ragionamento è possibile riporta
direttamente alla scelta di codifica descritta a inizio sezione: perché
la heatmap possa calcolare $P(A)$, $P(B)$ e $P(A,B)$ per una coppia di
tratti iconografici, è necessario poter contare, su tutte le schede del
corpus, quante volte compare ciascuna chiave e quante volte una coppia
di chiavi compare insieme --- un'operazione elementare se il dato è
strutturato come nel Codice~\ref{lst:xenodata}, impraticabile se fosse
rimasto testo libero. Il Codice~\ref{lst:pmi-pseudo} riassume, in forma
di pseudocodice, il collegamento tra il blocco \texttt{xenoData} di
ciascuna scheda e il calcolo della PMI descritto sopra: la struttura
dati intermedia non è altro che una tabella di conteggi, popolata
scandendo una sola volta il corpus.

\begin{lstlisting}[caption={Da \texttt{xenoData} alla heatmap: schema
del calcolo (pseudocodice)}, label={lst:pmi-pseudo}]
per ogni scheda S del corpus:
    tratti(S) = { (trait.type, trait.key)
                  per ogni trait in xenoData(S) }
    per ogni coppia non ordinata (A, B) in tratti(S):
        conteggio[A]      += 1
        conteggio[B]      += 1
        conteggio[(A, B)] += 1

N = numero totale di schede

per ogni coppia (A, B) con conteggio[(A, B)] > 0:
    P_A  = conteggio[A]      / N
    P_B  = conteggio[B]      / N
    P_AB = conteggio[(A, B)] / N
    PMI(A, B) = log2( P_AB / (P_A * P_B) )
\end{lstlisting}

Questo è, in ultima analisi, il senso più concreto della
\virg{comunicabilità} tra la sezione iconografica e la heatmap: non una
semplice possibilità in linea di principio, ma una conseguenza diretta e
quasi automatica dell'aver scelto, fin dall'inizio, chiavi controllate
al posto di prosa libera. Lo stesso identico meccanismo, con la stessa
tabella di conteggi, alimenta anche gli assi onomastici e geografici
della heatmap menzionati più sopra: l'iconografia non è un caso
particolare gestito a parte, ma un ulteriore asse dello stesso
strumento di analisi.

\screenshotplaceholder{7cm}{Sezione iconografica di una scheda, con i descrittori strutturati applicati}

\screenshotplaceholder{7.5cm}{Heatmap delle co-occorrenze tra attributi iconografici, onomastici e geografici}

\subsection{Onomastica ed epiteti divini}

Se la scheda descrittiva resta lo strumento naturale per consultare una
singola epigrafe, un corpus come quello raccolto da ILA pone anche
un'esigenza diversa: quella di orientarsi nell'insieme delle divinità e
delle figure onomastiche attestate, prima ancora di sapere quale singola
epigrafe si vuole consultare. A questo scopo sono state progettate due
viste di navigazione dedicate, concettualmente distinte pur condividendo
lo stesso principio ispiratore: restituire, in forma sintetica e
scorrevole, un indice che nella tradizione a stampa richiederebbe
sfogliare l'intero repertorio.

La prima vista riguarda le divinità e i rispettivi epiteti cultuali, e ne
propone una presentazione in forma di elenco alfabetico disposto lungo
una linea diagonale che attraversa la pagina, con le voci più recenti
progressivamente sfalsate verso destra --- una soluzione visiva debitrice
del linguaggio grafico dei manifesti o delle locandine espositive,
piuttosto che di quello, più severo, dei repertori scientifici
tradizionali. Ogni voce, selezionata mediante scorrimento continuo della
pagina, apre l'elenco delle epigrafi in cui quel nome, o quel
particolare epiteto, è attestato. La seconda vista riguarda invece il
registro onomastico più ampio (nomi propri, patronimici, cariche),
organizzato secondo il modello, familiare a chi ha consuetudine con
repertori e indici editoriali a stampa, della rubrica alfabetica
verticale, corredata da un indice rapido A--Z a margine della pagina che
consente di raggiungere direttamente una lettera senza dover scorrere
l'intero elenco.

Entrambe le viste dispongono di un campo di ricerca libera che filtra
l'elenco in tempo reale mentre l'utente digita. La ricerca è stata resa
volutamente sensibile a maiuscole e minuscole: nel dominio onomastico e
teonimico trattato dal corpus, iniziale maiuscola e minuscola non sono
equivalenti dal punto di vista filologico allo stesso modo in cui lo
sono, ad esempio, in una ricerca su prosa corrente, e restituire risultati
insensibili alla capitalizzazione avrebbe rischiato di confondere, in
alcuni casi, forme che nel testo originale sono tenute volutamente
distinte.

Nella vista diagonale, ogni divinità selezionata --- da ricerca o da
scorrimento --- apre un'anteprima ad albero dei propri epiteti, così da
mostrare la struttura interna del culto (quanti e quali epiteti sono
attestati per quel nome) senza dover aprire singolarmente ciascuna voce.
Quando il termine cercato individua un solo risultato possibile --- il
caso tipico è un epiteto che compare, nel corpus, associato a una sola
divinità --- quella voce viene automaticamente fissata come attiva e la
relativa anteprima ad albero si aggiorna di conseguenza, invece di
lasciare che sia l'utente a scorrere manualmente l'elenco per
raggiungerla: un piccolo accorgimento, ma coerente con l'idea generale
che guida l'intera sezione, ossia che orientarsi nell'onomastica del
corpus debba costare all'utente il minor numero di passaggi possibile.

\screenshotplaceholder{8cm}{Elenco diagonale delle divinità ed epiteti, in stile locandina}

\screenshotplaceholder{8cm}{Rubrica onomastica alfabetica con indice laterale A--Z}

\subsection{Le fonti letterarie e il lessico cultuale}

Il catalogo epigrafico racconta che cosa veniva inciso sulla pietra, ma
il culto di Men è documentato anche indirettamente, attraverso quanto ne
scrivevano gli autori greci e latini --- geografi, poeti, storici,
lessicografi. Per raccogliere anche questo secondo tipo di testimonianza
si è affiancata al catalogo una sezione di pari rango, dedicata alla
raccolta ragionata delle fonti letterarie relative alla divinità,
concepita fin dall'inizio non come una semplice appendice bibliografica
ma come una seconda collezione strutturata, interrogabile con gli stessi
criteri già impiegati per il materiale epigrafico.

L'impianto di questa sezione riprende, adattandolo, il modello di un
progetto di riferimento nello studio del rapporto tra lingua e religione
nell'antichità, dal quale si sono mutuate tre scelte di fondo: trattare
ogni citazione --- letteraria, epigrafica o papirologica --- come
un'unità autonoma e marcata in quanto tale; organizzare il materiale
attraverso un insieme di indici trasversali (le fonti, gli ambiti
tematici, i termini significativi, le divinità, i personaggi, le figure e
i luoghi menzionati); e classificare ciascuna testimonianza secondo una
griglia concettuale che incrocia alcuni campi d'analisi con altrettanti
ambiti d'uso del testo. A questo impianto si è aggiunto un elemento che
il progetto di riferimento non prevede: un collegamento diretto con il
catalogo epigrafico, per cui ogni testimonianza letteraria porta con sé
ricerche già pronte sul corpus delle epigrafi, espresse come
interrogazioni e non come riferimenti fissi a singole schede, in modo che
restino valide anche quando il catalogo continua a crescere.

La sezione è stata popolata a partire da una voce pilota dedicata a una
divinità lunare strettamente connessa a Men nella tradizione classica,
che raccoglie un consistente numero di testimonianze suddivise in nuclei
tematici disposti in ordine cronologico, dalla poesia arcaica greca fino
alla storiografia tardoantica, ciascuna corredata di testo originale,
traduzione redazionale e commento. L'ultimo di questi nuclei tematici è
dedicato specificamente al punto di contatto tra questa figura divina e
Men, e costituisce il ponte esplicito verso il corpus epigrafico.
Coerentemente con il principio, già più volte richiamato in questo
lavoro, per cui nessuna informazione viene presentata come acquisita
prima di essere stata verificata, ogni testo antico importato in questa
sezione è dichiaratamente segnalato come da collazionare sulle edizioni
di riferimento indicate, e l'interfaccia lo dichiara esplicitamente prima
ancora di mostrarne il contenuto.

Dal punto di vista della navigazione, questa sezione è stata inizialmente
concepita come un saggio continuo, da leggere dall'inizio alla fine;
questa forma, pur avendo il pregio della leggibilità, seguiva una logica
difforme da quella di tutto il resto della base dati, e non offriva un
livello di indice sopra la singola voce. È stata pertanto riorganizzata
su tre livelli, speculari a quelli già impiegati per il catalogo
epigrafico: un indice delle voci disponibili, che dichiara anche quelle
ancora da redigere; un elenco filtrabile delle singole testimonianze, con
la stessa impaginazione tabellare già impiegata per l'elenco delle
epigrafi, ordinabile per nucleo tematico, per cronologia o per autore; e
una scheda di dettaglio per la singola testimonianza, articolata in
sezioni per il testo, il commento, l'analisi e gli indici, i rimandi e la
bibliografia, sullo stesso modello della scheda epigrafica. La forma
discorsiva originaria non è stata eliminata, ma è diventata una modalità
di lettura ulteriore, accessibile in alternativa all'elenco e all'indice.

Un aspetto rilevante di questa sezione è che il testo letterario viene
marcato con lo stesso motore di codifica impiegato per l'edizione
epigrafica (si veda la sezione 3), del quale riusa le operazioni
concettualmente comuni ai due domini --- divinità, personaggi, luoghi,
lessico cultuale --- aggiungendo soltanto ciò che un testo letterario
richiede e una pietra incisa no (l'apparato critico, le parole citate in
quanto tali, i titoli d'opera, la distinzione tra tradizione letteraria,
epigrafica e papirologica), ed escludendo per contro ciò che appartiene
specificamente al supporto materiale dell'iscrizione. È proprio questa
base condivisa a rendere possibile il collegamento con il resto del
progetto: una divinità o un epiteto individuati all'interno di una
testimonianza letteraria confluiscono automaticamente negli indici di
questa sezione e compaiono, con rimando diretto, nelle pagine di
navigazione onomastica e divina già descritte nella sottosezione 5.2,
così che chi consulta la pagina di un epiteto possa raggiungere in un
clic le fonti letterarie che lo menzionano --- mentre il conteggio delle
attestazioni epigrafiche resta rigorosamente distinto, poiché risponde a
una domanda diversa (quanti monumenti, non quanti passi testuali).

Ogni testimonianza porta inoltre due indicazioni interpretative distinte,
anch'esse mutuate dal progetto di riferimento: una qualificazione
sintetica della natura della fonte --- se descrive una pratica cultuale
realmente osservata, una credenza, oppure un trattamento puramente
letterario o mitografico della divinità, dato che l'osservazione di un
geografo e una cosmogonia poetica non hanno lo stesso peso probatorio ---
e un'analisi più granulare applicata a singoli segmenti del testo.
Proprio perché lo scarto cronologico tra il momento di composizione di un
testo e l'epoca dei fatti o delle realtà che esso descrive resta, in
generale, un problema aperto e non sempre risolvibile in maniera
univoca, le fonti letterarie sono per il momento tenute fuori dalla
visualizzazione cronologica descritta in apertura di questa sezione, in
attesa di una regola esplicita e documentata per collocarle correttamente
lungo l'asse del tempo.

A questa sezione se ne affianca una seconda, più circoscritta ma
concettualmente contigua, dedicata al lessico cultuale: un insieme chiuso
di termini --- non nomi propri o epiteti, ma parole comuni legate a uno
specifico atto di culto, come un voto, una dedica, una colpa e la sua
espiazione, o una formula fissa ricorrente --- marcato con lo stesso
meccanismo sia nel testo delle epigrafi sia in quello delle fonti
letterarie. Poiché i due tipi di fonte condividono lo stesso vocabolario
controllato e la stessa marcatura, il medesimo lemma può essere osservato
attraverso entrambe, senza tuttavia confonderle in un unico totale che ne
offuscherebbe la diversa natura probatoria: ogni voce del lessico riporta
due conteggi distinti, uno per le epigrafi e uno per le fonti letterarie,
mai sommati fra loro, e i lemmi attestati soltanto nei testi letterari,
senza alcun riscontro epigrafico, sono segnalati come tali anziché essere
presentati con un'attestazione fittizia pari a zero. La visualizzazione
dispone i lemmi come una serie di barre, raggruppabili per famiglia
tematica o elencabili singolarmente per frequenza complessiva, la cui
lunghezza non è proporzionale al numero di attestazioni ma alla sua
radice quadrata:

\[
  \ell(A) \;\propto\; \sqrt{n(A)}
\]

dove $n(A)$ indica il numero di attestazioni del lemma $A$. Questa scelta
ha una motivazione puramente percettiva: alcune famiglie lessicali sono,
per la natura stessa del materiale, molto più frequenti di altre, e una
scala lineare le renderebbe leggibili a scapito delle famiglie più rare,
che si ridurrebbero a barre praticamente invisibili; comprimendo la scala
secondo la radice quadrata, le famiglie rare restano percepibili senza
che quelle dominanti perdano la propria posizione relativa. Ogni voce
della lista rimanda inoltre a uno strumento lessicografico greco esterno
per la consultazione del lemma, e consente di raggiungere direttamente le
singole attestazioni, epigrafiche o letterarie, da cui il conteggio
deriva.

\screenshotplaceholder{8cm}{Indice delle voci della sezione dedicata alle fonti letterarie}

\screenshotplaceholder{7.5cm}{Vista a barre del lessico cultuale, con attestazioni epigrafiche e letterarie a confronto}

Il lessico cultuale appena descritto dice a quale famiglia semantica
appartiene un termine --- un voto, una colpa, una formula di espiazione
--- ma lascia aperta una domanda distinta, altrettanto rilevante per uno
studio comparativo del culto: entro quale funzione più ampia si colloca,
nell'economia complessiva della pratica religiosa, l'atto o la
condizione che quel termine designa. Per rispondere a questa seconda
domanda si è integrata nel sistema una tassonomia funzionale di
riferimento, elaborata da un progetto dedicato allo studio comparato del
lessico delle iscrizioni confessionali anatoliche, articolata in un
ristretto numero di categorie principali --- l'espiazione, la
trasgressione, la condizione della persona coinvolta --- a loro volta
suddivise in sottocategorie più specifiche, innestate su due livelli
successivi.

Questa seconda classificazione non è stata trattata come una scelta
editoriale ulteriore da compiere voce per voce, ma come un'informazione
derivata dal lemma stesso, secondo lo stesso principio già enunciato a
proposito della rappresentazione dei dati (sezione 2): chi marca il
lessico cultuale continua a rispondere a una sola domanda --- quale
lemma ricorre in quel punto del testo --- e il collocamento all'interno
della tassonomia funzionale viene ricavato automaticamente da questa
scelta, anziché richiedere una seconda decisione indipendente che
rischierebbe, con il tempo, di disallinearsi dalla prima. Per la stessa
ragione la tassonomia non viene scritta stabilmente all'interno di
ciascuna epigrafe del corpus, ma calcolata al momento della
consultazione: la si scrive esplicitamente soltanto quando il documento
filologico è destinato a lasciare il progetto per essere condiviso con
chi cura la tassonomia di riferimento, così che quella collocazione
compaia dove effettivamente serve senza costituire, all'interno del
corpus stesso, una seconda fonte di verità da mantenere sincronizzata a
ogni ritocco della griglia esterna.

A questa nuova classificazione è stata affiancata una vista di
navigazione dedicata, che riprende la stessa resa grafica già impiegata
per il lessico cultuale --- barre in scala di radice quadrata, conteggio
a fianco --- disponendo però i lemmi secondo la struttura della
tassonomia funzionale anziché per famiglia semantica o per frequenza
assoluta. La vista mostra l'intera griglia, comprese le categorie e
sottocategorie per cui il corpus non offre, al momento, alcuna
attestazione: anche un ramo vuoto è un'informazione, poiché indica dove
la documentazione epigrafica e letteraria tace, e nasconderlo
restituirebbe un'immagine incompleta della tassonomia di riferimento. Le
voci proprie del progetto e quelle introdotte da un ampliamento della
griglia originaria restano inoltre segnalate distintamente, cosicché chi
consulta la vista possa sempre distinguere ciò che appartiene alla
tassonomia di riferimento da ciò che il progetto vi ha aggiunto in
proprio. Poiché la stessa suddivisione funzionale viene applicata sia al
lessico attestato nelle epigrafi sia a quello rilevato nelle fonti
letterarie, la griglia diventa un secondo punto di incontro tra le due
collezioni, distinto da quello già descritto a proposito delle divinità
e degli epiteti (sottosezione 5.2): ogni ramo della tassonomia riporta,
affiancate e mai sommate, le attestazioni provenienti da entrambe le
fonti.

\screenshotplaceholder{7.5cm}{Vista a griglia della tassonomia funzionale del lessico cultuale, con rami non attestati dichiarati come tali}

\subsection{Ricerca}

L'interfaccia digitale dispone di un modulo di ricerca che consente di
effettuare interrogazioni sulle epigrafi sui diversi attributi che le
caratterizzano: materiale del supporto, tipologia di monumento, luogo di
provenienza, lingua dell'iscrizione, datazione e, come illustrato nella
sottosezione precedente, singoli descrittori iconografici. È inoltre
disponibile una modalità di ricerca testuale libera sul testo delle
epigrafi, con normalizzazione dei segni diacritici greci, così da
restituire risultati coerenti indipendentemente dalla presenza o assenza
di accenti e spiriti nella stringa di ricerca inserita dall'utente.

I risultati di ciascuna operazione di ricerca vengono presentati in
forma di elenco, con accesso diretto alla scheda descrittiva completa di
ciascuna epigrafe individuata, in modo da rendere il passaggio dalla
ricerca alla consultazione del singolo monumento quanto più diretto e
intuitivo possibile.

Un caso d'uso frequente, soprattutto in fase di verifica o di
riferimento incrociato con altre schede, è la ricerca diretta di
un'epigrafe a partire dal suo numero identificativo. Poiché nel parlato
--- e altrettanto spesso nella scrittura corrente, specie quando ci si
riferisce a una voce a memoria --- il numero viene naturalmente
espresso in lettere anziché in cifre, il campo di ricerca riconosce
anche questa forma, convertendo internamente i numerali scritti per
esteso (\virg{novantuno}) nella corrispondente sequenza numerica
(\virg{91}) prima di interrogare l'indice: una piccola concessione
all'uso linguistico ordinario, che evita all'utente il passaggio mentale
da parola a cifra che il solo campo numerico avrebbe altrimenti
richiesto.

\screenshotplaceholder{7cm}{Interfaccia del modulo di ricerca con filtri combinati}

La ricerca appena descritta risponde a un'esigenza di tipo
catalografico: individuare le epigrafi che condividono un determinato
attributo. Chi lavora invece direttamente sul testo delle iscrizioni ha
spesso una domanda diversa, di natura più propriamente filologica e
linguistica: dove ricorre una determinata parola nel corpus, e che cosa
la circonda. Per rispondere a questa seconda esigenza si è introdotta
una sezione di concordanza, che restituisce ciascuna occorrenza del
termine cercato accompagnata dal proprio contesto immediato, secondo il
modello consolidato della concordanza a parola-chiave-in-contesto
(\textit{key word in context}, KWIC), ben noto a chi si occupa di
linguistica dei corpora.

La concordanza legge lo stesso testo già impiegato dal modulo di ricerca
generale, senza richiedere un indice dedicato: la sola differenza è che,
in questo caso, l'unità restituita è la singola occorrenza in contesto e
non la scheda nel suo complesso. Anche qui le parole spezzate da
un'interruzione di riga --- un fenomeno frequente nel testo epigrafico
trascritto secondo le convenzioni descritte nella sezione 3 --- vengono
ricongiunte prima della ricerca, cosicché un termine che nel testo
originale attraversa un cambio di riga resti comunque cercabile nella
sua forma intera; e la ricerca resta, come altrove nel modulo di ricerca
generale, insensibile a differenze di accentazione e di maiuscole. I
risultati possono essere ordinati per il contesto che precede o che
segue l'occorrenza --- un'operazione che, applicata a un insieme
sufficientemente ampio di attestazioni, fa emergere da sola le formule
ricorrenti, senza che sia necessario cercarle esplicitamente ---
filtrati per regione di provenienza, ed esportati in un formato
tabellare per un'elaborazione al di fuori dell'interfaccia.

\screenshotplaceholder{7cm}{Sezione di concordanza, con le occorrenze disposte in contesto}

\section{Considerazioni finali e sviluppi futuri}

Il lavoro descritto in questo articolo rappresenta un primo bilancio
delle attività di sviluppo del progetto ILA, condotte parallelamente su
più fronti: il riconoscimento ottico e la trascrizione dei testi a
partire dall'edizione a stampa di riferimento, l'adozione e
l'adattamento di un sistema di marcatura filologica consolidato, la
revisione sistematica della classificazione di divinità ed epiteti
sull'intero corpus, e la progettazione di un'interfaccia di consultazione
che affianca, agli strumenti tradizionali di ricerca epigrafica, una
sezione dedicata alla descrizione strutturata dell'iconografia dei
monumenti --- aspetto che, come si è detto, risultava finora trattato
solo in forma discorsiva nella letteratura di riferimento --- e due viste
di navigazione pensate specificamente per orientarsi nell'insieme delle
divinità, degli epiteti e delle figure onomastiche attestate nel corpus,
affiancate da una visualizzazione cronologica delle attestazioni e da una
sezione, di pari rango rispetto al catalogo epigrafico, dedicata alla
raccolta delle fonti letterarie relative alla divinità e al lessico
cultuale che le accomuna al materiale inciso sulla pietra. A queste
attività si sono affiancati, più di recente, un primo strumento di
collazione assistita tra il corpus e le edizioni a stampa di
riferimento, una tassonomia funzionale che affianca al lessico cultuale
una seconda chiave di lettura, una sezione di concordanza pensata per
chi lavora direttamente sul testo delle iscrizioni, e un apparato di
citazione stabile per ciascuna voce del corpus.

L'approccio complessivamente seguito privilegia la trasparenza e la
verificabilità del dato rispetto alla comodità di gestione a breve
termine: si è scelto di non introdurre alcun formato intermedio
semplificato tra il dato grezzo e la sua codifica filologica, e di non
integrare mai automaticamente informazioni non direttamente attestate
dalle fonti. Questa scelta comporta un onere maggiore nelle fasi di
sviluppo e di revisione, ma garantisce che ogni informazione presentata
all'utente sia sempre riconducibile, in maniera diretta e verificabile,
alla fonte da cui deriva.

Per quanto riguarda gli sviluppi futuri, il lavoro prosegue
principalmente su due fronti. Da un lato si sta portando avanti la
revisione sistematica, voce per voce, delle trascrizioni derivate
dall'edizione a stampa originaria, sulla base di un confronto diretto
con le fonti primarie disponibili --- confronto per cui si dispone ora
di un primo strumento di collazione assistita, descritto nella sezione
4, ma che resta da estendere sistematicamente all'intero corpus --- così
da ridurre progressivamente la dipendenza dal testo di Lane e
trasformarlo, per ciascuna voce revisionata, da testo sottostante a
semplice riferimento bibliografico.
Dall'altro lato si prevede di ampliare ulteriormente il vocabolario
controllato utilizzato per la descrizione iconografica, e di
sviluppare strumenti di confronto comparativo tra le diverse
attestazioni iconografiche del dio Men, in modo da valorizzare appieno
il potenziale interrogativo di questa nuova sezione dell'interfaccia
digitale. Si prevede infine di estendere la sezione dedicata alle fonti
letterarie oltre la voce pilota già redatta, applicando lo stesso
modello alle altre figure divine e alle altre tradizioni testuali
rilevanti per il culto di Men, e di stabilire una regola esplicita e
documentata per la collocazione delle testimonianze letterarie
all'interno della visualizzazione cronologica, dalla quale restano per
ora escluse in ragione dello scarto, difficilmente risolvibile in via
generale, tra il momento di composizione di un testo e l'epoca dei fatti
che esso racconta.

\end{document}
